\chapter{Listado de código}
\label{ap:codigo}

Se reproduce aquí el fragmento de código que constituye la contribución
computacional central del trabajo: la derivada analítica exacta. El script
\texttt{py/generar\_extractos\_codigo.py} permite extraerlo del módulo de
valoración. El listado corresponde a la versión utilizada al generarlo; si el
código cambia, debe regenerarse y contrastarse con la revisión que produjo los
resultados del \cref{cap:resultados}. Este procedimiento reduce la
transcripción manual, sin garantizar por sí solo la sincronización de versiones
(véase el \cref{ap:software}).

El resto del código --- los \datoNMetodos{} métodos del catálogo y su
instrumentación, el estimador del orden de convergencia, la construcción del
problema, el protocolo de comparación y la recursión de retiro programado ---
forma parte del proyecto. El repositorio indicado es \repositorio; para
reproducir los resultados se requieren además la revisión del código y las
entradas de la ejecución, cuya disponibilidad no se deduce del enlace.

\section{La derivada analítica exacta}
\label{sec:cod-derivada}

Este es el aporte central sobre el motor de valoración: la implementación de las
dos derivadas laterales de \cref{eq:subdif-der,eq:subdif-izq}, es decir, del
$B$-subdiferencial del \cref{thm:subdiferencial}. El cálculo reutiliza los
vectores de valoración y añade máscaras y productos escalares. El oráculo conjunto
devuelve valor y derivadas en una llamada al motor, evitando las valoraciones
auxiliares de las diferencias finitas; su costo marginal se mide junto con la
preparación y las comprobaciones de cada método.

Obsérvese en particular la corrección de la rama de crecimiento en $M=S_0$
(\cref{prop:continuidad-rama}), sin la cual la derivada en ese punto resulta
errónea; el factor de $34$ corresponde al ejemplo del \cref{cap:estructura}.

\lstinputlisting[caption={Derivadas laterales exactas de la reserva matemática
(fragmento de \texttt{core/motor.py}).}, label={lst:derivada}]%
{codigo/derivada_exacta.py}

% ---------------------------------------------------------------------
%  Listados trasladados al repositorio público (31-oct-2026). Se conservan
%  aquí, excluidos de la compilación, por si se quieren reincorporar:
%  basta con borrar las dos líneas que los encierran (iffalse ... fi).
% ---------------------------------------------------------------------
\iffalse

\section{Instrumentación del catálogo}

Todos los métodos comparten la misma instrumentación, sin lo cual la comparación
no sería justa. Dos decisiones merecen comentario.

La primera es que cada ejecución dispone de memoria propia: repetir una consulta
ya disponible no implica una nueva llamada al motor, pero la información no se
hereda de otro método. La contabilidad incluye preparación y derivadas; el oráculo
conjunto evita duplicar la valoración cuando entrega también las pendientes.

La segunda es que el criterio de parada es \emph{mixto}, en residuo y en paso, por
las razones discutidas en la \cref{crit:parada}.

\lstinputlisting[caption={Contador de evaluaciones y registro de iterados
(fragmento de \texttt{core/metodos\_iterativos.py}).}, label={lst:instrumentacion}]%
{codigo/instrumentacion.py}

\section{Estimación del orden de convergencia}

El estimador clásico de tres términos consecutivos es inservible aquí, por las
razones expuestas en la \cref{rem:estimacion-orden}. Se implementa en su lugar una
regresión de mínimos cuadrados sobre la cola útil de la sucesión de errores, junto
con dos estimadores auxiliares --- la razón lineal y la razón de reducción del ancho
del intervalo --- que son los descriptores adecuados para los métodos de orden
uno.\footnote{La razón de reducción del ancho del intervalo tiene además la virtud
de no requerir el conocimiento de la raíz. Cuando se dispone de cotas de
pendiente, el residuo proporciona otro certificado de error.}

\lstinputlisting[caption={Estimadores robustos del orden y de la razón de
convergencia.}, label={lst:orden}]{codigo/estimacion_orden.py}

\section{Métodos con encaje}

\subsection{Bisección}

Se incluye como referencia por su cota geométrica de error
(\cref{thm:cota-biseccion}), que no exige derivabilidad. Otros métodos con
encaje también ofrecen garantías bajo sus respectivas condiciones; la no
suavidad no hace exclusiva la garantía de bisección.

\lstinputlisting[caption={Método de bisección.}, label={lst:biseccion}]%
{codigo/metodo_biseccion.py}

\subsection{Regula falsi y sus tres correcciones}

Las cuatro variantes comparten una única implementación parametrizada por el
factor $\gamma$ con que se reduce el valor de $f$ en el extremo estancado. Esto no
es solo economía de código: hace explícito que Illinois, Pegasus y
Anderson--Björck difieren \emph{únicamente} en esa elección, que es exactamente lo
que la \cref{prop:estancamiento} y su discusión sostienen.

\lstinputlisting[caption={Regula falsi con las variantes de Illinois, Pegasus y
Anderson--Björck.}, label={lst:regula}]{codigo/metodo_regula_falsi.py}

\subsection{ITP}

El método más reciente del catálogo \citep{oliveira2021}, con sus tres etapas de
interpolación, truncamiento y proyección claramente separadas en el código.

\lstinputlisting[caption={Método ITP (\emph{Interpolate--Truncate--Project}).},
label={lst:itp}]{codigo/metodo_itp.py}

\section{Métodos abiertos}

\subsection{Secante}

\lstinputlisting[caption={Método de la secante, con acotamiento al rango
admisible.}, label={lst:secante}]{codigo/metodo_secante.py}

\subsection{Newton: las dos variantes}

El listado contiene las dos implementaciones que el \cref{cap:resultados}
contrasta: la derivada por diferencias centradas, con valoraciones auxiliares, y
la derivada exacta obtenida junto con la reserva. Los contadores incluyen las
llamadas efectivas y la preparación; la comparación de costo y precisión se
presenta en el \cref{cap:resultados}, sin atribuir un porcentaje fijo al número
nominal de consultas por iteración.

\lstinputlisting[caption={Newton con diferencias finitas y Newton con derivada
analítica.}, label={lst:newton}]{codigo/metodo_newton.py}

\section{Métodos para el caso no suave}

El primero dispone de la garantía local del \cref{cor:qi-sun-aplica}; el segundo
incorpora encaje y bisección periódica para la garantía global del
\cref{thm:salvaguardado}. Estas propiedades sustentan su análisis sin excluir
otros métodos que satisfagan sus propias condiciones.

\lstinputlisting[caption={Newton semisuave de Qi--Sun y su variante con
salvaguarda de encaje.}, label={lst:semisuave}]{codigo/metodo_semisuave.py}

\section{Construcción del problema}

\subsection{Función objetivo y subdiferencial}

Nótese el cambio de convención de signo respecto del motor de producción, que
define la función objetivo como $\KA-\VP(M)$, decreciente. Aquí se usa
$f=\VP-\KA$, creciente, que es la convención bajo la cual la teoría de métodos
para funciones monótonas crecientes se aplica sin cambios de signo.

\lstinputlisting[caption={Construcción de $f$, $\VP'_+$ y $\VP'_-$.},
label={lst:objetivo}]{codigo/objetivo.py}

\subsection{Capital mínimo y detección numérica de los quiebres}

Las dos funciones que verifican computacionalmente los resultados estructurales
del \cref{cap:estructura}: el capital mínimo financiable del
\cref{cor:capital-minimo} y la localización de los puntos de quiebre predichos por
el \cref{thm:quiebres}.

\lstinputlisting[caption={Verificación numérica de la estructura.},
label={lst:capital}]{codigo/capital_minimo.py}

\section{Protocolo de comparación}

Aquí queda fijado, en código, el protocolo descrito en el
\cref{cap:resultados}: mismo intervalo de encaje, mismo punto inicial, mismo
criterio de parada, misma contabilidad de costo para los \datoNMetodos{}
métodos. La preparación de parámetros, incluida una consulta de pendiente para
punto fijo, pertenece a su costo y no se suministra como información gratuita.

\lstinputlisting[caption={Ejecución del catálogo completo sobre un escenario.},
label={lst:banco}]{codigo/banco_pruebas.py}

\section{Recursión de retiro programado}

El bucle que implementa el sistema dinámico de la \cref{def:sistema-rp}. Se
reconocen los ajustes de la \cref{sec:operadores}, con techo comercial opcional y
desactivado por defecto, y la regla de arranque igual a la renta vitalicia. El
estado incluye la mesada anterior. La parada de la simulación por insuficiencia
de capital exige distinguir el cierre numérico del derecho pensional.

\lstinputlisting[caption={Bucle principal de la proyección de retiro programado
(fragmento de \texttt{core/retiro\_programado.py}).}, label={lst:rp}]%
{codigo/retiro_programado.py}

\fi
